\documentclass[10pt]{article}

\usepackage[utf8]{inputenc}

\usepackage[T1]{fontenc}

\usepackage{amsmath}

\usepackage{amsfonts}

\usepackage{amssymb}

\usepackage[version=4]{mhchem}

\usepackage{stmaryrd}

\usepackage{mathrsfs}

\usepackage{bbold}



\begin{document}

вень соответствия не столь высок либо из-за менышей точности эксперимента, либо вследствие существенного влияния эффектов, выходящих за рамки чистой К. э.



Вообще опытные данные по всем без исключения эффектам К.э. находятся в прекрасном согласии с теоретич. значениями в тех случаях, когда влияние других видов взаимодействий оказывается несушественным либо поддается учету. Этот факт имеет принципиальное значение как для К.э., так и для квантовой теории поля в целом. Он свидетельствует о том, что основные положения современной локальной (калибровочной) квантовой теории поля, а также динамич. основа К. э., соответствующая калибровочному лагранжиану взаимодействия, справедливы во всей области, доступной современному эксперименту.



Лит.: [1] А хи ез е р А.И., Бе р с тецки й В.Б., Квантовая электродинамика, 4 изд., М., 1981; [2] Бе рестецкий В. Б., Л и фш и ц Е. М., П и т а в ск и й Л. П., Релятивистская квантовая электродинамика, 2 изд., М., 1980; [3] Бо го л ю бо в Н. Н., Ш и р к о в Д. В., Введение в теорию квантованных полей, 4 изд., М., 1984; [4] Ф е й н м а н Р., Квантовая электродинамика, пер. с англ., М., 1964. Д. И. Казаков, Д. В. Ширков. KBAHTOBA9 OHTPOпИЯ - понятие, обобщающее понятие энтропии Больцмана и Гиббса, введенное Дж. Нейманом (J. Neumann, 1927). Формула



\[

S(\omega)=-k \operatorname{Tr}\left(\rho_{\omega} \log \rho_{\omega}\right)

\]



задает энтропию состояния $\omega$ с матрицей плотности $\rho_{\omega}$. B формализме Неймана ограниченным наблюдаемым физич. системы отвечают самосопряженные элементы $A$ алгебры $\mathscr{B}(H)$ всех ограниченных линейных операторов в гильбертовом пространстве $H$ и $\omega(A)=\operatorname{Tr}\left(\rho_{\omega} A\right)-$ среднее значение $A$ в состоянии $\omega$. Если



\[

\rho_{\omega}=\sum_{n} p_{n}\left|p_{n}\right\rangle\left\langle p_{n}\right|, \quad p_{n} \geqslant 0

\]



\begin{itemize}

  \item разложение по собственным функциям оператора плотности $\rho_{\omega} \in \mathscr{B}(H)$ со следом $\operatorname{Tr} \rho_{\omega}=\sum_{n} p_{n}=1$, то (в системе единиц, где постоянная Больцмана $h=1$ )

\end{itemize}



\[

S(\omega)=-\sum_{n} p_{n} \log p_{n} \in[0,+\infty]

\]



где $0 \cdot \log 0=0$. К. э. характеризует информацию, заключенную в состоянии $\omega$. К. э. не является наблюдаемой, это функционал на множестве состояний системы. Среди основных свойств К. э. каК математич. объекта (неотрицательность, вогнутость, полунепрерывность и др.) важна для приложений сильная субаддитивность, установленная позже других при помощи нек-рого свойства вогнутости и связана с исследованием квантовой относительной энтропии. Можно аксиоматически задать К. э. ее свойствами.



В квантовой статистич. физике, изучающей обладающие термодинамич. поведением бесконечные системы, применимость (1) ограничивается тем, что алгебры $\mathscr{B}(H)$ с $\operatorname{dim} H \leqslant \infty$ используются для описания конечных подсистем. Для широкого класса моделей по аналогии с классич. статистич. механикой можно ввести плотность энтропии (среднюю энтропию) трансляционно инвариантного состояния $\omega$ как нек-рый предел



\[

s(\omega)=\lim _{V \rightarrow \infty} S\left(\rho_{V} /|V|\right)<\infty,

\]



где матрица плотности $\rho_{V}$ отвечает сужению $\omega$ на подсистему в конечном объеме $V$.



В математич. теории классич. энтропии, в эргодич. теории классич. динамич. систем (с дискретным временем) важную роль играют инвариант Колмогорова-Синая, метрическая (динамическая) энтропия. Есть разные подходы к построению ее квантового аналога. Так, напр., можно опираться на свойства относительной К. э., тем же методом вводится энтропия для состояний на гиперфинитных $C^{*}-$ и $W^{*}$-алгебрах, применяемых в квантовой статистич. механике.



А. В. Булинский. КВАНТОВОЕ БРОУНОВСКОЕ ДВИХЕНИЕ - см. Квантовое стохастическое исчисление.



KВАНТОВОЕ ГОМОЛОГИЧЕСКОЕ УРАВНЕНИЕ уравнение для фазы $C$ усредняющего оператора:



\[

[H, C]=L-\underline{L}

\]



где $H$ - невозмущенный гамильтониан, $L-$ возмущающий оператор, $L-$ секулярный гамильтониан. Если $H$ является функцией от квантовых операторов действия и отсутствуют резонансы, то уравнение (1) разрешимо (см. [1], [2]) и в классич. пределе $(h \rightarrow 0)$ уравнение (1) переходит в гомологич. уравнение Пуанкаре (см. [3]).



Jum.: [1] We in ste in A., "Duke Math. J.», 1977, v. 44, p. 883-92; [2] Карасев М. В., Маслов В. П., «Успехи математич. наук», 1984, т. 39, в. 6, с. 115-73; [3] Арн ольд В. И., Дополнительные главы теории обыкновенных дифференциальных уравнений, М., $1978 . \quad$ Ю. М. Воробьев.



КВАНТОВОЕ ЕВКЛИДОВО ПОЛЕ - см. Евклидово квантовое поле.



KBAНTOBOE ИЗМЕРЕНИЕ - процесс получения данных, несущих информацию о состоянии квантовой системы, в основе которого лежит физическое взаимодействие системы с измерительным прибором. Характерная черта К. и. состоит в принципиальной неустранимости воздействия измерительного прибора на наблюдаемую систему, в результате чего ее состояние изменяется в ходе измерения.



Начала теории К. и. были заложены в монографии [1] Дж. Неймана (1932). В 1960-70-е гг. появление таких дисциплин, как квантовая оптика, оптич. связь, квантовая электроника, а также внутренние потребности квантовой теории вызвали новый интерес к вопросам измерений и привели к созданию последовательной статистич. теории К. и., являющейся логич. развитием вероятностной интерпретации квантовой механики и свободной от нек-рых трудностей и противоречий старого подхода (см. [4], [5], [7]).



Математическое описание квантовых измерений. Пусть $H$ - гильбертово пространство квантовой системы, $X-$ множество возможных исходов данного К. и. («шкала» измерительного прибора). Как правило, можно считать, что $X$ - стандартное измеримое пространство с $\sigma$-алгеброй измеримых подмножеств $B(X)$. Описание К. и. сушественно зависит от уровня подробности задания процесса измерения. Имеется три основных уровня, каждому из к-рых отвечает определенный математич. объект в гильбертовом пространстве $H$.



\begin{enumerate}

  \item Задана только статистика исходов измерения. Это означает, что для любого состояния системы, описываемого плотности оператором $\rho$, задано распределение вероятностей $\mu_{\rho}(B) ; B \in \mathscr{B}(X)$, исходов данного измерения над системой, приготовленной в этом состоянии $\left(\mu_{\rho}(B)\right.$ есть вероятность того, что исход измерения $x$ попадает в подмножество $B \subset X$ ). Естественно считать, что отображение $\rho \rightarrow \mu_{\rho}$ обладает свойством аффинности: смесь операторов плотности $\sum_{j} p_{j} \rho_{j}$ переходит в такую же смесь $\sum_{j} p_{j} \mu_{\rho_{j}}$ соответствующих распределений. При этом условии имеет место формула

\end{enumerate}



\[

\mu_{\rho}(B)=\operatorname{Tr} \rho \mathscr{\mu}(B), B \in \mathscr{B}(X),

\]



где $\mathscr{\mu}(\cdot)$ есть нек-рое разложение единицы в $H$. Обратно, всякое разложение единицы определяет по формуле (1) аффинное отображение $\rho \rightarrow \mu_{\rho}$ множества операторов плотности в множество распределений вероятностей на $X$. Разложение единицы дает наиболее общее описание статистики исходов измерения, совместимое с вероятностной интерпретацией квантовой механики.



П р име ры. a) Пусть $X-$ наблюдаемая, а $E(B)$, $B \in \mathscr{B}(\mathbb{R}),-$ ее спектральная мера на вещественной прямой $\mathbb{R}$. Формула (1) при $\mu(B) \equiv E(B)$ дает распределение вероятностей исходов точного измерения наблюдаемой $X$ в состоянии $\rho$.



б) Пусть $|x, v\rangle$ есть вектор состояния наименьшей неопределенности (см. Волновой пакет) для нерелятивистской квантовой частицы массы $m$ со средними значениями коор-





\end{document}